\documentclass[10pt,twocolumn]{article}

\usepackage[margin=0.72in]{geometry}
\usepackage[T1]{fontenc}
\usepackage[utf8]{inputenc}
\usepackage{microtype}
\usepackage{graphicx}
\usepackage{subcaption}
\usepackage{booktabs}
\usepackage{array}
\usepackage{tabularx}
\usepackage[hidelinks]{hyperref}
\usepackage{xcolor}
\usepackage{enumitem}
\usepackage{amsmath}

\graphicspath{{../../../runs/memory_policy/20260525_114946/}}

\setlength{\columnsep}{0.22in}
\setlength{\parindent}{0pt}
\setlength{\parskip}{0.45em}
\renewcommand{\arraystretch}{1.12}
\setlist[itemize]{leftmargin=1.2em,itemsep=0.2em,topsep=0.25em}
\makeatletter
\setlength{\@fptop}{0pt}
\setlength{\@fpbot}{0pt plus 1fil}
\setlength{\@fpsep}{8pt}
\makeatother

\definecolor{PaperBlue}{HTML}{203A63}
\definecolor{PaperOrange}{HTML}{F59E0B}
\definecolor{PaperGray}{HTML}{5B6578}
\definecolor{PaperLight}{HTML}{F8FAFC}

\title{\vspace{-1.1em}\bfseries Five Memory Policies in a 5M Transformer: A Mac Smoke Sweep}
\author{Vuk Rosic\\\small Independent Research Project}
\date{May 25, 2026}

\begin{document}
\maketitle

\begin{abstract}
We compare six attention styles in the same small transformer scaffold:
dense attention, age-based forgetting, usage-refresh memory, competition-based
memory, hierarchical summarization, and predictive importance memory. The goal
here is not to claim a final winner. It is to verify that the memory
philosophies can be implemented, trained, and instrumented in one shared code
path on a MacBook, while still producing readable plots and tables for later
NVIDIA runs. This PDF is a smoke report, so its 2048-token training budget is
only for the debug pass; the intended research sweep uses 5,000,000 tokens per
condition and includes a plain CSA baseline as a seventh comparison point.
In this smoke sweep, validation loss is essentially tied across the memory
policies, while dense attention is much faster in plain PyTorch.
\end{abstract}

\section{What We Are Testing}

The research question is simple:

\begin{quote}
If we keep the model and data fixed, do five different memory philosophies
produce meaningfully different learning dynamics?
\end{quote}

The five policies are:

\begin{itemize}
\item age-based forgetting,
\item usage-refresh memory,
\item competition-based memory,
\item hierarchical summarization memory,
\item predictive importance memory.
\end{itemize}

Dense attention is the baseline. It gives every token access to the full causal
past. The memory policies instead keep a local window and route older context
through compressed or scored memory blocks.

This report is a smoke test, not the final paper claim. The real purpose is to
show that the code path is stable, the outputs are legible, and the same
framework can later be moved to a longer NVIDIA run.

\section{Setup}

Table~\ref{tab:setup} gives the shared configuration. The model is the
compact 5M-parameter preset discussed earlier in the repo, with a 256-token
context window. For this Mac smoke sweep, we only train for 2048 tokens so the
run finishes quickly and the plots are easy to inspect. The real research
budget is 5,000,000 tokens per condition.

\begin{table}[t]
\centering
\small
\caption{Shared setup for the Mac smoke sweep. The 2048-token budget here is
only for the smoke pass; the real research sweep uses 5,000,000 tokens per
condition.}
\label{tab:setup}
\begin{tabular}{@{}ll@{}}
\toprule
Item & Value \\
\midrule
Model preset & 5M-parameter transformer \\
Hidden size & 256 \\
Layers / heads & 6 / 8 \\
Context window & 256 \\
Batch size & 4 \\
Smoke budget & 2048 tokens \\
Target research budget & 5,000,000 tokens/condition \\
Seeds & 1 (Mac smoke) \\
Dataset & synthetic causal copy-lag \\
Hardware & MacBook / MPS debug run \\
Main metric & validation loss \\
\bottomrule
\end{tabular}
\end{table}

The fairness rule is:

\begin{quote}
Same model, same data, same batch size, same context length, same optimizer
settings, and the same tiny training budget.
\end{quote}

The only thing that changes is attention policy.

For the memory policies, we also track cheap post-run probe metrics. That keeps
the training loop lightweight and pushes the extra instrumentation into a
separate pass.

\section{Results}

Table~\ref{tab:results} summarizes the sweep.

\begin{table*}[t]
\centering
\scriptsize
\setlength{\tabcolsep}{3pt}
\caption{Mac smoke results. Peak CUDA memory is not meaningful on this MPS/CPU debug run, so the table emphasizes loss, throughput, and the probe metrics.}
\label{tab:results}
\begin{tabular}{@{}lrrrrrrrrrr@{}}
\toprule
Policy & Loss & Tok/s & Vec. & Cov. & Gate & Refresh & Util. & Pred. & Sel. & Lvl. \\
\midrule
dense & 8.3776 & 1809.5 & 256 & 256 & -- & -- & -- & -- & -- & -- \\
age-based forgetting & 8.3772 & 166.9 & 80 & 128 & 0.3792 & -- & -- & -- & 1.88 & -- \\
usage-refresh & 8.3772 & 161.6 & 80 & 128 & 0.4204 & 0.1250 & -- & -- & 1.88 & -- \\
competition & 8.3772 & 159.7 & 80 & 128 & -- & -- & -2.4697 & -- & 1.44 & -- \\
hierarchical & 8.3772 & 159.0 & 84 & 128 & 0.9167 & -- & -- & -- & 2.25 & 2.0 \\
predictive & 8.3798 & 148.6 & 80 & 128 & -- & -- & -- & 0.5001 & 1.44 & -- \\
\bottomrule
\end{tabular}
\end{table*}

The main visual result is that the validation losses are almost indistinguish-
able on this tiny run, while throughput separates the dense baseline from the
memory policies very clearly. To make that easier to read, we use a single
summary figure with loss deltas and throughput ratios instead of two nearly
identical line charts.

\begin{figure*}[t]
\centering
\includegraphics[width=0.96\textwidth]{policy_comparison_summary.png}
\caption{Policy comparison summary. Top-left: absolute validation loss, zoomed
to the narrow range where the policies differ. Top-right: loss delta relative
to dense, in milli-loss. Bottom-left: throughput on a log scale. Bottom-right:
throughput relative to dense. This makes the near-ties in loss and the large
throughput gap much easier to read at a glance.}
\label{fig:summary}
\end{figure*}

\begin{figure*}[t]
\centering
\includegraphics[width=0.98\textwidth]{memory_mechanisms_overview.png}
\caption{One-panel schematics for the five memory philosophies. Each panel
shows the qualitative mechanism rather than the exact kernel implementation:
age decay, refresh, competition, hierarchy, and predictive retention.}
\label{fig:mechanisms}
\end{figure*}

\begin{figure}[t]
\centering
\includegraphics[width=\linewidth]{gate_by_age.png}
\caption{Age-based forgetting gate strength decreases with memory age. This is
the one policy with a visibly monotonic retention curve in the smoke run.}
\label{fig:gate}
\end{figure}

\section{Interpretation}

The smoke run says three useful things.

\begin{itemize}
\item The shared code path works. All six policies train and emit metrics.
\item The memory policies are visually distinct. The probe plots are not the
same shape, which is exactly what we wanted from the new abstractions.
\item This tiny run is not enough to claim a real quality win. Dense is still
the throughput leader, and the validation losses are almost tied.
\end{itemize}

That is a good result for a tutorial pipeline. It means the experiment is
usable, not yet persuasive.

\section{What This Does Not Prove}

This run does \emph{not} prove that any memory philosophy is better than dense.
It only proves that we can compare them in a consistent harness and capture the
behavior in plots and tables.

The next real step is the same comparison on NVIDIA, with more tokens and a few
seeds. Then the question becomes interesting in the research sense:

\begin{quote}
Do the five memory philosophies separate once the model has enough data and
enough runtime?
\end{quote}

\section{Conclusion}

We now have a paper-ready path for the new memory work:

\begin{quote}
shared model $\rightarrow$ five memory policies $\rightarrow$ one fairness
rule $\rightarrow$ automated plots and tables $\rightarrow$ PDF report.
\end{quote}

On the Mac smoke run, the policies are distinguishable and the reporting is in
place, but the losses are still too close to support a substantive claim. That
is exactly the kind of honest result a mini research paper should report.

\begin{thebibliography}{9}
\bibitem{deepseekv4}
DeepSeek-AI.
\newblock {\em DeepSeek-V4: Towards Highly Efficient Million-Token Context Intelligence}.
\newblock 2026.

\bibitem{pytorch}
A.~Paszke et al.
\newblock PyTorch: An Imperative Style, High-Performance Deep Learning Library.
\newblock 2019.
\end{thebibliography}

\end{document}
